\documentclass[10pt,a4paper]{amsart}
\usepackage[margin=19mm]{geometry}
\usepackage{amsmath,amssymb,mathtools}
\usepackage[scr=rsfs,cal=euler]{mathalfa}
\usepackage{xfrac}
\usepackage[unicode,hidelinks,bookmarks=false]{hyperref}
\newcommand{\ie}{i.e.\ }
\newcommand{\cf}{cf.\ }
\newcommand{\resp}{respectively\ }
\newcommand{\Subsection}{\subsection}
\providecommand{\nobreakdash}{\nobreak\hbox{-}}
\setlength{\emergencystretch}{2em}
\multlinegap0pt
\usepackage[shortlabels]{enumitem}
\usepackage{amssymb}
\newtheorem{lemma}{Lemma}
\newcommand{\Cs}{\mathrm{C}^*}
\newcommand{\Csr}{\mathrm{C}^*_{\text{red}}}
\title{Higher Kazhdan projections and delocalised $\ell^ 2$-Betti numbers}
\author{Sanaz Pooya, Hang Wang}
\date{Source: arXiv:2405.03837v2 (CC BY 4.0)}
\begin{document}
\maketitle

\noindent\emph{Source and licence.} Higher Kazhdan projections and delocalised $\ell^ 2$-Betti numbers. Version arXiv:2405.03837v2 (CC BY 4.0), distributed by arXiv under the Creative Commons Attribution 4.0 International licence, http://creativecommons.org/licenses/by/4.0/, as recorded in the arXiv licence metadata for this exact version and shown on the abstract page https://arxiv.org/abs/2405.03837v2. Full licence notice: \texttt{licenses/arxiv-2405.03837-CC-BY-4.0.txt}.\par\smallskip

\noindent\textit{Editorial scope.} Counted unit: the article's lemma lem:projker (source0.tex lines 1263--1281). The article's complete original proof of it is delivered with it (source0.tex lines 1283--1377, one \texttt{proof} environment of 95 lines). No mathematics is written, repaired or re-derived: the statement and the proof are the article's own lines, in the article's own notation, and no retained line is substituted or re-flowed.  Retained uncounted as the source-local context the proof needs: lines 1092--1106; lines 1237--1241; lines 1244--1247.  Nothing was omitted from any retained span, so the package carries no omission record.  No retained line contains a citation, so the wrapper carries no bibliography and no bibliography entry.  No retained line contains a figure environment, an image-inclusion command or drawing code, so no image is delivered and the package declares no asset dependency.  Editorial formatting only, in addition: an AMS \texttt{amsart} preamble that restates the article's own declarations and macros used by the retained lines, namely the environments \texttt{lemma} on separate counters, together with the standard packages those lines need.  The retained environments print in this document as Lemma 1 in order of appearance, whereas the article prints the same items with the numbering of its own full document.  The complete native source of the article as distributed in the arXiv e-print for this exact version is bundled unchanged as \texttt{sources/158755/source0.tex}.
	\begin{lemma}
		\label{lem:SolDelta1}
		Let $G= \mathbb Z_m \ast \mathbb Z_n$ with $m \geq 2$ and $n \geq 3$. 
		Let $x\in \ell^2 (G)^{\oplus 2}$. Then 
		$\Delta_1 x =0$ 
		if and only if there are 
		$z\in \ell^2 (G)^{\oplus 2}$ and 
		$a\in \mathrm{im}(1-p)\cap\mathrm{im}(1-q)$
		such that 
		\[
		  x=(I - C)z  \qquad \text{and} \qquad		
		  \begin{bmatrix}1-s & 0 \\ 0 & {1-t}\end{bmatrix}z= \frac{1}{\sqrt 2}\begin{bmatrix}a\\ -a\end{bmatrix}. 
		\] 
		% for some $a\in \mathrm{im}(1-p)\cap\mathrm{im}(1-q).$
	\end{lemma}

	\begin{equation}
		\label{eq3}
		1-p = \frac{1}{m}[(1-s)+\cdots+(1-s^{m-1})] =
		k(1-s)=(1-s)k,
	\end{equation}

	\begin{equation}
		\label{eq4}
		1-q = l(1-t)=(1-t)l.
	\end{equation}

	\begin{lemma}
		\label{lem:projker}
		Let $k$ and $l$ be as in (\ref{eq3}) and (\ref{eq4}). The $\ell^2$-kernel of $\Delta_1$ is   
		\[
		\ker\Delta_1
		=
		\left\{\begin{bmatrix}ka\\ -la\end{bmatrix} \mid  a\in \mathrm{im}(1-p)\cap\mathrm{im}(1-q)\right\}.
		\] 
		Moreover, the projection $P_{\ker \Delta_1}$ has the same $K$-class as the projection $P_{H_1}$ onto the Hilbert subspace 
		\begin{equation}
			\label{eq:H1}
			H_1:=\left\{\begin{bmatrix}a\\ -a\end{bmatrix} \mid  a\in \mathrm{im}(1-p)\cap\mathrm{im}(1-q)\right\}
			\subset \ell^2(G)\oplus \ell^2(G).
		\end{equation}
		In other words,  
		\[
			[P_{\ker\Delta_1}]=[P_{H_1}]\in \mathrm K_0(\Csr(G)).
		\]
	\end{lemma}

	\begin{proof}
		%Let $H=l^2(G).$
		In view of (\ref{eq3}), (\ref{eq4}) and Lemma \ref{lem:SolDelta1}, we have that $\Delta_1 x=0$ if and only if $x$ satisfies  
		\[
			x
			=
			\begin{bmatrix}1-p & 0 \\ 0 & 1-q\end{bmatrix}z
			= 
			\begin{bmatrix}k & 0 \\ 0 & l\end{bmatrix}
			\begin{bmatrix}1-s & 0 \\ 0 & 1-t \end{bmatrix}
			z
			=
			\begin{bmatrix}k & 0 \\ 0 & l\end{bmatrix}
			\begin{bmatrix}\frac{1}{\sqrt 2}a \\-\frac{1}{\sqrt 2}a
			\end{bmatrix} 
		\]
		for $z=(I-C)x$ and some $a\in \ell^2(G)$ satisfying $a\in\mathrm{im}(1-p)\cap\mathrm{im}(1-q)$.
		In other words, we have
		\begin{align*}
			\ker\Delta_1
			=
			&
			\left\{ 	
				\begin{bmatrix}1-p & 0 \\ 0 & 1-q\end{bmatrix}z
				~~\mid~~  
				\begin{bmatrix}1-s & 0 
				\\ 
				0 & 1-t\end{bmatrix}z=\begin{bmatrix}\frac{1}{\sqrt 2}a \\-\frac{1}{\sqrt 2}a
			\end{bmatrix}, a\in\mathrm{im}(1-p)\cap\mathrm{im}(1-q)  	
		   \right\}
		   \\
			=
			&
			  \left\{ 
			  	\begin{bmatrix}\frac{ka}{\sqrt 2} \\ \frac{-la}{\sqrt 2}\end{bmatrix}
			 	 ~~\mid~~ 
			  	a\in\mathrm{im}(1-p)\cap\mathrm{im}(1-q)  
			  \right\}.
		\end{align*}
	%	Due to the factorization in (\ref{eq3}) and (\ref{eq4}), $k$ and $1-s$ are inverse to each other in $(1-p)\ell^2(G)$ and $l$ and $1-t$ are so in $(1-q)\ell^2(G)$. 
	%With respect to the decomposition 
	%$\ell^2(G) = p\ell^2(G)\oplus(1-p)\ell^2(G)$, and together with the fact that 
	%$1-s=(1-p)(1-s)$ and $1-t=(1-q)(1-t)$
	%the operator $v_1:=p+(1-s)$	is invertible with the inverse given by $v_1^{-1}:= 1+k$.
	%The same holds for the operator $v_2:=q+(1-t)$. Its inverse is given by $v_2^{-1}:= 1+l$.
Due to the factorizations in (\ref{eq3}) and (\ref{eq4}), the operators $k$ and $1-s$ act as mutual inverses on the subspace $(1-p) \ell^2(G)$, just as $l$ and $1-t$ act as mutual inverses on $(1-q) \ell^2(G)$. To see this, note that $k(1-p)=(1-p)k$, and since $1-s=(1-p)(1-s)$, the space $(1-p)\ell^2(G)$ is invariant under the action by $1-s$. With respect to the orthogonal decomposition $\ell^2(G)=p \ell^2(G) \oplus(1-p) \ell^2(G)$, the operator $v_1:=p+(1-s)$ respects this block-diagonal structure. It acts as the identity on $p \ell^2(G)$ and as $1-s$ on $(1-p) \ell^2(G)$. Therefore, $v_1$ is invertible with its inverse given by $v_1^{-1}:= p + k(1-p)$. 
Indeed, a direct calculation verifies this (noting that $sp = p$):
\begin{align*}
[p+(1-s)][p+k(1-p)] &= p^2 + pk(1-p) + (1-s)p + (1-s)k(1-p) \\
&= p + 0 + (p-sp) + (1-p) \\
&= p + 0 + 0 + (1-p) \\
&= 1.
\end{align*}


Consider the invertible operator $V:=\begin{bmatrix} v_1 & 0 \\ 0 & v_2 \end{bmatrix}$ in $M_2(\Csr(G))$. Let us compute the image of $\ker \Delta_1$ under $V$. For $a \in \mathrm{im}(1-p)$, we have $pa = 0$. This implies that $pka = pk(1-p)a = 0$. Consequently, the action of $v_1$ on $ka$ yields:
\[
v_1(ka) = (p+1-s)ka = pka + (1-s)ka = 0 + (1-p)a = a.
\]
Applying the same logic to the second coordinate, we find:
\begin{align*}
V(\ker \Delta_1) &= \left\{ \begin{bmatrix} v_1 ka \\ -v_2 la \end{bmatrix} \;\middle|\; a \in \mathrm{im}(1-p) \cap \mathrm{im}(1-q) \right\} \\
&= \left\{ \begin{bmatrix} a \\ -a \end{bmatrix} \;\middle|\; a \in \mathrm{im}(1-p) \cap \mathrm{im}(1-q) \right\} = H_1.
\end{align*}

To conclude the statement in K-theory, consider the element $e := V^{-1} P_{H_1} V \in M_2(\Csr(G))$. Since $P_{H_1}$ is an orthogonal projection onto $H_1$, $e$ is an idempotent operator whose image is precisely $V^{-1}(H_1) = \ker \Delta_1$. In a $\Cs$-algebra, any idempotent defines the same $\mathrm K_0$-class as the orthogonal projection onto its image, so $[P_{\ker \Delta_1}] = [e] \in \mathrm K_0(\Csr(G))$. Furthermore, since $e$ and $P_{H_1}$ are conjugate idempotents, they share the same K-theory class. Thus, we obtain:
\[
[P_{\ker \Delta_1}] = [e] = [V^{-1} P_{H_1} V] = [P_{H_1}].
\]
This finishes the proof.
		% originaly writen by Hang
		%We then have
		%\begin{equation*}
			%V(\ker \Delta_1)
			%=
		%	\ker(V \Delta_1 V^{-1}) 
		%	\quad \text{and} \quad 
		%	P_{\ker\Delta_1}
		%	=
		%	V^{-1}P_{\ker(V\Delta_1 V^{-1})}V.
		%\end{equation*}
		%This implies that as $K$-theory elements
		%\[
		%[P_{ker\Delta_1}]=[P_{\ker(V\Delta_1V^{-1})}],
		%\]
	%	where the kernel of ${(V\Delta_1V^{-1})}$ satisfies
		%\[
		%\ker(V\Delta_1 V^{-1})
		%=
		%V(\ker \Delta_1) 
		%= 
	%\]
		%which is isomorphic to $\left\{\begin{bmatrix}a\\ -a\end{bmatrix} \mid  a\in \mathrm{im}(1-p)\cap\mathrm{im}(1-q)\right\}$
		%as desired.
	\end{proof}

\end{document}
